% English translation of prelude.tex, lines 185--205.
% Source: Methods of Algebra, Volume 2, commit
% 9a5803ff77dd3257484cb177f851a73770a59dd3 (CC BY 4.0).
% stable-unit-id: o014.aljabr2.prelude.acknowledgments

% segment-id: o014.li2.prelude.acknowledgments.h001
\section*{Acknowledgments}
\addcontentsline{toc}{section}{Acknowledgments}
\hypertarget{o014.aljabr2.prelude.acknowledgments}{ }

% segment-id: o014.li2.prelude.acknowledgments.p001
During the preparation of this book, the author benefited from corrections and suggestions by teachers, colleagues, and friends. The process lasted so long that omissions and errors of memory are inevitable; the following list is incomplete: Huan Zhen, Huang Chenxin, Li Jinghui, Lei Jiale, Li Changyuan, Sun Yuze, Tang Yiming, Xue Hanyu, Yang Enlin, Zhang Haofeng, and Zhao Zelong (in Hanyu Pinyin order).

% segment-id: o014.li2.prelude.acknowledgments.p002
A small portion of the book was used in courses at the University of Chinese Academy of Sciences and Huazhong University of Science and Technology. The author also thanks everyone who took part in those courses.

% segment-id: o014.li2.prelude.acknowledgments.p003
Like all the author's works to date, this book was written entirely in an open-source software environment, including the operating system and the fonts used in preparing it for press. The author salutes the many selfless developers who put into practice the universal values of sharing and altruism.

% segment-id: o014.li2.prelude.acknowledgments.p004
While this book was being written, the author's research was supported by the Excellent Young Scientists Fund of the National Natural Science Foundation of China under project 11922101, with Peking University as the host institution. This support is gratefully acknowledged.

% segment-id: o014.li2.prelude.acknowledgments.p005
The author also thanks Zhao Tianfu, an editor at Higher Education Press, for his patience and professionalism. The writing of this book ought to have followed seamlessly from Volume~I, but various obstacles delayed its completion until the end of 2022.

% segment-id: o014.li2.prelude.acknowledgments.p006
Despite the delay, preparing the book consumed even more of the author's mental energy than Volume~I. This was especially true of the repeated checking and revision in the final stages, which felt like walking through an endless dark tunnel. Much of this late work was done in Beijing subway cars. Working on the train was hardly comfortable and required making every minute and second count, quite literally. Yet amid that immense, anonymous current, day after day, the author seemed to sense a kind of tacit mutual understanding rarely found in working life, and to draw a glimmer of consolation from it. Was it so, or was it not? The explanation may be fanciful, but the feeling was real. The prelude now draws to a close: homage to the boundless distance, and to the endless multitudes.

% segment-id: o014.li2.prelude.acknowledgments.sig001
\vspace{1em}
\begin{flushright}\begin{minipage}{0.3 \textwidth}
	\begin{tabular}{c}
		% segment-id: o014.li2.prelude.acknowledgments.sig002
		{\fangsong Wen-Wei Li} \\
		% segment-id: o014.li2.prelude.acknowledgments.sig003
		December 2022 \\
		% segment-id: o014.li2.prelude.acknowledgments.sig004
		Shijingshan
	\end{tabular}
\end{minipage}\end{flushright}
